\begin{explanationbox}
	\textbf{\textcolor{CExplanation}{一句话直觉}}

	从圆外一点引切线与割线，切线长的平方等于割线上两段线段长的乘积。

	\medskip
	\textbf{\textcolor{CExplanation}{核心拆解}}
	\begin{description}[style=nextline, leftmargin=2.8em, labelsep=0.5em, itemsep=0.45em]
		\item[\textbf{要点一（相似构造）：}] 连接切点与割线两端点，构造两个共点相似三角形 $\triangle PAT$ 与 $\triangle PTB$。
		\item[\textbf{要点二（等积变形）：}] 由相似得比例式 $\dfrac{PA}{PT}=\dfrac{PT}{PB}$，再转化为等积式 $PT^{2}=PA\cdot PB$。
	\end{description}

	\medskip
	\textbf{\textcolor{CExplanation}{几何本质（三角相似）}}

	在 $\triangle PAT$ 与 $\triangle PTB$ 中，由弦切角定理得 $\angle PTA=\angle PBT$，又 $\angle APT=\angle BPT$（公共角），故 $\triangle PAT\sim\triangle PTB$，从而得到比例式。

	\medskip
	\textbf{\textcolor{CExplanation}{代数意义（比例中项）}}

	等式 $PT^{2}=PA\cdot PB$ 表明 $PT$ 是 $PA$ 与 $PB$ 的比例中项，体现了线段之间的等量关系。

	\medskip
	\textbf{\textcolor{CExplanation}{考点价值（常考应用）}}
	\begin{itemize}[leftmargin=*, itemsep=0.5em, topsep=0.4em]
		\item \textbf{考法一（求切线长）：} 已知割线两段长度，直接代入公式求切线长。
		\item \textbf{考法二（求割线线段）：} 已知切线长和一段割线长，求另一段。
	\end{itemize}

	\medskip
	\textbf{\textcolor{CExplanation}{顿悟点}}

	关键条件：切线、割线必须从同一点出发，且 $A$ 在 $P$、$B$ 之间。

	\medskip
	\textbf{\textcolor{CExplanation}{使用场景}}
	\begin{itemize}[leftmargin=*, itemsep=0.5em, topsep=0.4em]
		\item \textbf{场景一（切线割线共点）：} 圆外一点与圆的一条切线和一条割线同时出现时。
		\item \textbf{场景二（圆幂计算）：} 涉及圆外点到圆的幂时，可利用 $d^{2}-R^{2}$ 直接计算。
	\end{itemize}
\end{explanationbox}
